\scanpage{19}\mspage{18}
$L$ est muni d'une topologie, comme celle d'être \uline{localement} fini. Le sous-ens.
ordonné ${\unsure{\mathcal F|\mathscr F}}'$ \deleted{de} $\mathcal F$ est alors stable par sup \unsure{fini|finis}, mais pas \deleted{par} \unsure{généralt} par sup.
quelc. \uline{Figure finie} $F$ : \deleted{finie} $F$ fini au sens ordinaire, i.e. un \unsure{ens|nb} fini de strates. Figure de type fini : contenue
dans une réunion finie de ss-figures $F_A$.

\deleted{Si les strates d'une} Une figure est \deleted{\unsure{de t.}} finie \deleted{si} si et seult si elle est de t. f. et ses strates sont finies.

Autre cas : supposons $L$ muni d'une structure d'ordre -- il \unsure{--|a} dans une réalisation topologique, par des ``\uline{cellules}'' \inserted{(fermées) $X_x$ $(x\in L)$} dans un espace
top. $X$ triangulé convenable. Alors à tout $A\deleted{\unsure{\in}}\unsure{\subset|\in}L$ correspond \deleted{\unsure{le}} $X_A=\bigcup_{x\in A}X_x$.
% slanted
\margin{les $X_x$ forment une \uline{famille} admissible de parties de $X$, indexée par l'ens ordonné $L$}
On se bornera alors alors aux parties \uline{fermées} $A$ de $L$, de façon que
\[ A\longmapsto X_A \]
soit une application bijective de l'ens des parties envisagées de $L$ (les p. \uline{fermées})
\deleted{\unsure{dans l'}}\inserted{\unsure{avec un|sur un}} ens des parties \inserted{fermées} de $X$ \deleted{\unsure{stables}} \deleted{par la relation} \inserted{stables par la} relation d'équivalence \unsure{induite|induit} par \deleted{\unsure{les}} $(X_x^\circ)_{x\in L}$ (en fait, toutes les parties
